\documentclass[12pt,a4paper]{article}
\usepackage[utf8]{inputenc}
\usepackage[T1]{fontenc}
\usepackage[spanish]{babel}
\usepackage{graphicx}
\usepackage{booktabs}
\usepackage{hyperref}
\usepackage[backend=biber,style=ieee]{biblatex}
\addbibresource{referencias.bib}

\begin{document}

% ---------- Carátula ----------
\begin{titlepage}
  \centering
  {\Large Universidad Técnica Estatal de Quevedo\par}
  {\large Facultad de Ciencias de la Computación\par}
  {\large Carrera de Ingeniería de Software\par}
  \vspace{1cm}
  {\large Asignatura: Aplicaciones Distribuidas\par}
  {\large Actividad: TA-IND-01\par}
  \vspace{1cm}
  {\Large\bfseries Análisis crítico de un sistema distribuido real\par}
  \vspace{1cm}
  Estudiante: NOMBRE\par
  Paralelo: PARALELO\par
  Fecha: FECHA\par
  Repositorio: \url{https://github.com/ameraj2-bit/ta-ind-01-sistemas-distribuidos}
\end{titlepage}

% ---------- Resumen ----------
\begin{abstract}
% 150 a 200 palabras: sistema analizado, enfoque y conclusiones principales.
\end{abstract}

\tableofcontents
\newpage

\section{Introducción}
\section{Arquitectura del sistema}
\section{Análisis con los conceptos de la Unidad 1}
\subsection{Transparencias de distribución}
\subsection{CAP y PACELC}
\subsection{Modelos de comunicación}
\subsection{Tolerancia a fallos}
\section{Experimentación propia}
\section{Ejemplos del entorno del estudiante}
\section{Análisis comparativo}
\section{Postura crítica}
\section{Conclusiones}
\section{Declaración de uso de inteligencia artificial}

\printbibliography

\appendix
\section{Esquema inicial a mano}
\section{Capturas de lectura de las fuentes}

\end{document}
